\documentclass{article}
\usepackage[margin=1in]{geometry}
\usepackage{amsmath,amssymb,amsthm}
\usepackage{booktabs,longtable}
\usepackage{hyperref}

\newtheorem{definition}{Definition}
\newtheorem{theorem}[definition]{Theorem}
\newtheorem{lemma}[definition]{Lemma}
\newtheorem{proposition}[definition]{Proposition}
\newtheorem{convention}[definition]{Convention}
\newtheorem{warning}[definition]{Warning}

\newcommand{\source}[1]{\par\noindent\textbf{Source.} #1}
\newcommand{\lean}[1]{\par\noindent\textbf{Lean target.} \texttt{\detokenize{#1}}}
\newcommand{\status}[1]{\par\noindent\textbf{Status.} \texttt{\detokenize{#1}}}
\newcommand{\depends}[1]{\par\noindent\textbf{Depends on.} #1}
\newcommand{\cA}{\mathcal A}
\newcommand{\Ftwo}{\mathbb F_2}
\newcommand{\tr}{\operatorname{tr}}
\newcommand{\spanF}{\operatorname{span}_{\mathbb F_2}}

\begin{document}

\section*{Blueprint 009: Internal Stabilizer Reduction}

\subsection*{Scope}

This note promotes the second remaining blueprint block from
\path{agents/2026-06-22-remaining-blueprint-suite-goal.md}.  It supports
targets
\[
  \texttt{SYMP-FIN-001},\qquad
  \texttt{SYMP-COUNT-BASIS-001},\qquad
  \texttt{SYMP-FIN-RESTRICT-ISCOMPL-001},\qquad
  \texttt{SYMP-ORTH-RESTRICT-NONDEG-001},\qquad
  \texttt{SYMP-LININD-001},\qquad
  \texttt{SYMP-HYP-EXT-001},\qquad
  \texttt{SYMP-COUNT-FIN-001},\qquad
  \texttt{SYMP-PAULI-FIN-COUNT-001},\qquad
  \texttt{SYMP-INF-REC-001},\qquad
  \texttt{SYMP-SPANNING-BASIS-001},\qquad
  \texttt{SYMP-ADAPTED-CLOSED-001},\qquad
  \texttt{SYMP-INF-CE-001},\qquad
  \texttt{SYMP-CLOSED-PARTNER-001},\qquad
  \texttt{SYMP-SUP-HYP-NONDEG-001},\qquad
  \texttt{SYMP-ORTH-COMP-001},\qquad
  \texttt{SYMP-WSPLIT-PRESERVE-001},\qquad
  \texttt{SYMP-PAIR-RESIDUAL-001},\qquad
  \texttt{SYMP-STEP-CLOSED-001},\qquad
  \texttt{SYMP-W-EQ-SPAN-001},\qquad
  \texttt{SYMP-FRESH-PERP-001},\qquad
  \texttt{SYMP-BOT-NONDEG-001},\qquad
  \texttt{SYMP-SEQUENCES-CLOSED-001},\qquad
  \texttt{SYMP-INF-REC-CLOSED-001},\qquad
  \texttt{SYMP-001},\qquad
  \texttt{SYMP-002},\qquad
  \texttt{PROD-001}.
\]
It does not add Lean declarations, Aristotle packets, operator-algebra type
claims, or any local-circuit/gLU claim.  The goal is only the algebraic
stabilizer-side reduction used later by the GNS and type-classification
layers.

\subsection*{Source Basis}

The main project-local source is
\begin{center}
\path{../runs/2026-06-18-stabilizer-algebra-classification/artifacts/004-internal-stabilizer-reduction-type-classification-R1.pdf}.
\end{center}
The relevant source locators are:
\begin{itemize}
\item Lemma 3.1, PDF pp.~3--4: adapted symplectic basis for a countable
isotropic subspace.
\item Lemma 3.2, PDF p.~4: adapted symplectic basis induces a Pauli-algebra
\(*\)-isomorphism.
\item Proposition 4.1, PDF pp.~4--5: product form after symplectic basis
reduction.
\item Warning 6.2, PDF p.~7: the basis change is algebraic and need not be
geometrically local.
\end{itemize}
The finite Pauli and quasi-local algebra conventions are supplied by
\path{blueprints/007-pauli-quasilocal-algebra.tex}.  The region-restriction
and crossing-invisibility input is supplied by
\path{blueprints/008-stabilizer-local-state-foundations.tex}.

\subsection*{Conventions}

\begin{convention}[Region label space and internal stabilizers]
Let \(\Lambda\) be a finite or countably infinite region.  Its finite-support
binary Pauli label space is
\[
  V_\Lambda=(\Lambda\to_0\Ftwo)\oplus(\Lambda\to_0\Ftwo)
\]
with the binary symplectic commutation form \(\sigma\).  Let
\(W\subseteq V_\Lambda\) be an isotropic internal stabilizer subspace.  In the
application \(W=W'_\Lambda=W'\cap V_\Lambda\), where \(W'\) is the stabilizer
label space of an extension \(S'\).
\source{\path{blueprints/008-stabilizer-local-state-foundations.tex},
targets \texttt{PAULI-001}, \texttt{QL-002}, and \texttt{STATE-004}; project
proof note 004, Lemma 2.4, PDF p.~3.}
\status{Convention for all targets in this note.}
\end{convention}

\begin{convention}[Symplectic basis notation]
A symplectic basis of \(V_\Lambda\) is a family
\[
  \{z_i,x_i:i\in I\}
\]
such that every vector of \(V_\Lambda\) is a finite \(\Ftwo\)-linear
combination of the \(z_i,x_i\), and
\[
  \sigma(z_i,z_j)=0,\qquad
  \sigma(x_i,x_j)=0,\qquad
  \sigma(z_i,x_j)=\delta_{ij}.
\]
The index set \(I\) is finite or countable.
\source{Project proof note 004, Lemma 3.1, PDF pp.~3--4.}
\status{Notation for \texttt{SYMP-FIN-001}, \texttt{SYMP-001}, and
\texttt{SYMP-002}, and for the checked support targets used in the countable
assembly.}
\end{convention}

\subsection*{Adapted Symplectic Basis}

\begin{lemma}[FiniteAdaptedSymplecticBasis, ID: SYMP-FIN-001, project proof note 004 Lemma 3.1, PDF pp.~3--4]
Let \(\Lambda\) be finite, so that \(V_\Lambda\) is finite-dimensional over
\(\Ftwo\).  Let \(W\subseteq V_\Lambda\) be isotropic for the nondegenerate
symplectic form \(\sigma\).  Then there are a natural number \(n\), vectors
\[
  z,x:\{0,\ldots,n-1\}\longrightarrow V_\Lambda,
\]
and a subset \(P\subseteq\{0,\ldots,n-1\}\) such that \(z,x\) form a
symplectic basis of \(V_\Lambda\) and
\[
  W=\spanF\{z_i:i\in P\}.
\]
\source{Project proof note 004, Lemma 3.1, finite-dimensional part of the
proof, PDF pp.~3--4.}
\depends{\texttt{PAULI-001}.}
\lean{Formalization.Stabilizer.adaptedSymplecticBasis_finite}
\status{lean_checked supporting finite case for \texttt{SYMP-001}.}
\end{lemma}

\begin{proof}
The finite proof is the symplectic Gram--Schmidt construction from Lemma 3.1
with termination supplied by finite dimension.  If \(W=0\), choose a nonzero
vector \(z_0\), use nondegeneracy to find \(x_0\) with
\(\sigma(z_0,x_0)=1\), split off the hyperbolic plane
\(\spanF\{z_0,x_0\}\), and recurse on its orthogonal complement.  If
\(W\neq0\), choose \(0\neq z_0\in W\), find \(x_0\) with
\(\sigma(z_0,x_0)=1\), and use isotropy to decompose
\[
  W=\spanF\{z_0\}\oplus \bigl(W\cap \spanF\{z_0,x_0\}^{\perp}\bigr).
\]
The restriction of \(\sigma\) to the orthogonal complement is again
nondegenerate and alternating.  Finite-dimensional induction on the orthogonal
complement gives adapted vectors for the intersection term, and adjoining
\((z_0,x_0)\) gives the desired finite adapted symplectic basis.
\end{proof}

\subsection*{Checked Support Targets From the Countable Assembly}

\begin{lemma}[CountableSubmoduleBasis, ID: SYMP-COUNT-BASIS-001, project proof note 004 Lemma 3.1, PDF pp.~3--4]
Let \(V\) be a countable \(\Ftwo\)-vector space and let
\(W\subseteq V\) be a subspace.  Then there are a type \(S\), a basis
\[
  b:S\longrightarrow W
\]
of \(W\) as an \(\Ftwo\)-vector space, and a countability structure on \(S\).
\source{Project proof note 004, Lemma 3.1, countable-basis setup for the
countable recursive construction, PDF pp.~3--4; standard vector-space basis
existence.}
\depends{\texttt{PAULI-001}.}
\lean{Formalization.Stabilizer.exists_countable_basis_submodule}
\status{lean_checked support target returned by
\texttt{SYMP-001-COUNTABLE-ASSEMBLY-2026-06-24}.}
\end{lemma}

\begin{proof}
Choose a basis of \(W\) by the vector-space basis theorem.  Since \(W\) is a
subtype of the countable type \(V\), \(W\) is countable.  The basis map is
injective, so its index type injects into a countable type and is therefore
countable.
\end{proof}

\begin{lemma}[FiniteRestrictIsCompl, ID: SYMP-FIN-RESTRICT-ISCOMPL-001, project proof note 004 Lemma 3.1, PDF pp.~3--4]
Let \(V\) be an \(\Ftwo\)-vector space with an alternating bilinear form
\(B\).  Let \(H\subseteq V\) be finite-dimensional.  If the restriction
\(B|_H\) is nondegenerate, then \(H\) and its orthogonal complement
\[
  H^\perp=\{v\in V:B(v,h)=0\text{ for all }h\in H\}
\]
are complementary subspaces:
\[
  V=H\oplus H^\perp .
\]
No finite-dimensionality assumption on \(V\) is required.
\source{Project proof note 004, Lemma 3.1, finite-stage hyperbolic splitting
used inside the countable construction, PDF pp.~3--4.}
\depends{\texttt{SYMP-FIN-001}.}
\lean{Formalization.Stabilizer.finite_restrict_isCompl}
\status{lean_checked support target returned by
\texttt{SYMP-001-COUNTABLE-ASSEMBLY-2026-06-24}.}
\end{lemma}

\begin{proof}
The intersection \(H\cap H^\perp\) is zero by nondegeneracy of \(B|_H\).
For \(u\in V\), the functional \(h\mapsto B(h,u)\) on the finite-dimensional
space \(H\) is represented by a unique \(h_0\in H\), because \(B|_H\) is
nondegenerate.  Then \(u-h_0\) is orthogonal to \(H\), so
\(u=h_0+(u-h_0)\in H+H^\perp\).  Thus the sum is direct and exhaustive.
\end{proof}

\begin{lemma}[OrthogonalRestrictNondegenerate, ID: SYMP-ORTH-RESTRICT-NONDEG-001, project proof note 004 Lemma 3.1, PDF pp.~3--4]
Let \(B\) be a nondegenerate alternating bilinear form on \(V\).  Let
\(H\subseteq V\) be finite-dimensional and suppose \(B|_H\) is
nondegenerate.  Then the restriction of \(B\) to \(H^\perp\) is
nondegenerate.
\source{Project proof note 004, Lemma 3.1, recursive passage to the
orthogonal complement, PDF pp.~3--4.}
\depends{\texttt{SYMP-FIN-RESTRICT-ISCOMPL-001}.}
\lean{Formalization.Stabilizer.orthogonal_restrict_nondeg}
\status{lean_checked support target returned by
\texttt{SYMP-001-COUNTABLE-ASSEMBLY-2026-06-24}.}
\end{lemma}

\begin{proof}
By \texttt{SYMP-FIN-RESTRICT-ISCOMPL-001}, every \(v\in V\) decomposes
uniquely as \(h+h'\), with \(h\in H\) and \(h'\in H^\perp\).  If an element
of \(H^\perp\) pairs trivially with all of \(H^\perp\), it also pairs
trivially with \(H\) by orthogonality, hence with all of \(V\).  The
nondegeneracy of \(B\) on \(V\) then forces it to be zero.  The same argument
applies on the right side of the bilinear form.
\end{proof}

\begin{lemma}[SymplecticLinearIndependentGeneral, ID: SYMP-LININD-001, project proof note 004 Lemma 3.1, PDF pp.~3--4]
Let \(V\) be an \(\Ftwo\)-vector space, let \(B\) be an
\(\Ftwo\)-valued bilinear form on \(V\), let \(I\) be any index type, and let
\(z,x:I\to V\) satisfy
\[
  B(z_i,z_j)=0,\qquad B(x_i,x_j)=0,\qquad
  B(z_i,x_j)=\delta_{ij}
\]
for all \(i,j\in I\).  Then the combined family
\[
  I\oplus I\longrightarrow V,\qquad
  \operatorname{inl}(i)\mapsto z_i,\quad
  \operatorname{inr}(i)\mapsto x_i
\]
is linearly independent over \(\Ftwo\).
\source{Project proof note 004, Lemma 3.1, independence check for the
finite-support symplectic family in the countable construction, PDF
pp.~3--4.}
\depends{\texttt{SYMP-FIN-001}.}
\lean{Formalization.Stabilizer.symplectic_linearIndependent_general}
\status{lean_checked support target returned by
\texttt{SYMP-001-COUNTABLE-ASSEMBLY-2026-06-24}.}
\end{lemma}

\begin{proof}
A linear dependence has finite support.  Pairing the dependence with \(x_k\)
extracts the coefficient of \(z_k\), and pairing it with \(z_k\) extracts the
coefficient of \(x_k\).  All coefficients are therefore zero.
\end{proof}

\begin{lemma}[HyperbolicExtension, ID: SYMP-HYP-EXT-001, project proof note 004 Lemma 3.1, PDF pp.~3--4]
Let \(B\) be a nondegenerate alternating bilinear form on \(V\).  Let
\(H\subseteq V\) be finite-dimensional and suppose \(B|_H\) is
nondegenerate.  For any \(t\notin H\), there are \(p,q\in H^\perp\) such that
\[
  B(p,q)=1,\qquad p\neq0,\qquad
  t\in H+\spanF\{p\}.
\]
\source{Project proof note 004, Lemma 3.1, one-step hyperbolic extension in
the recursive adapted-basis construction, PDF pp.~3--4.}
\depends{\texttt{SYMP-FIN-RESTRICT-ISCOMPL-001},
\texttt{SYMP-ORTH-RESTRICT-NONDEG-001}.}
\lean{Formalization.Stabilizer.exists_hyperbolic_extension}
\status{lean_checked support target returned by
\texttt{SYMP-001-COUNTABLE-ASSEMBLY-2026-06-24}.}
\end{lemma}

\begin{proof}
Split \(t=h+p\) using \(V=H\oplus H^\perp\).  Since \(t\notin H\), the
orthogonal component \(p\) is nonzero.  By nondegeneracy of \(B\) on
\(H^\perp\), choose \(q\in H^\perp\) with \(B(p,q)=1\).  The decomposition
also gives \(t\in H+\spanF\{p\}\).
\end{proof}

\begin{theorem}[CountableAdaptedFiniteDimensional, ID: SYMP-COUNT-FIN-001, project proof note 004 Lemma 3.1, PDF pp.~3--4]
If \(V\) is finite-dimensional over \(\Ftwo\), \(B\) is nondegenerate and
alternating, and \(W\subseteq V\) is isotropic, then there are a countable
type \(I\), vectors \(z,x:I\to V\), and \(P\subseteq I\) such that
\[
  \{z_i,x_i:i\in I\}
\]
is a symplectic basis and \(W=\spanF\{z_i:i\in P\}\).
\source{Project proof note 004, Lemma 3.1, finite-dimensional part recast
with a countable index, PDF pp.~3--4.}
\depends{\texttt{SYMP-FIN-001}.}
\lean{Formalization.Stabilizer.exists_adapted_symplectic_countable_of_finiteDimensional}
\status{lean_checked support target returned by
\texttt{SYMP-001-COUNTABLE-ASSEMBLY-2026-06-24}.}
\end{theorem}

\begin{proof}
Apply the checked finite adapted-basis theorem and take \(I\) to be the
finite index type of the resulting basis.  Every finite type is countable.
\end{proof}

\begin{theorem}[FiniteRegisterCountableIndexAdapted, ID: SYMP-PAULI-FIN-COUNT-001, project proof note 004 Lemma 3.1, PDF pp.~3--4]
For a finite qubit label type \(\Gamma\), every isotropic subspace
\[
  W\subseteq V_\Gamma
\]
has an adapted symplectic basis with the same countable-index output shape as
\texttt{SYMP-001}.
\source{Project proof note 004, Lemma 3.1, finite-dimensional part recast
for \(\texttt{PauliLabel}\), PDF pp.~3--4.}
\depends{\texttt{PAULI-001}, \texttt{SYMP-COUNT-FIN-001}.}
\lean{Formalization.Stabilizer.adaptedSymplecticBasis_countableIndex_of_finite}
\status{lean_checked support target returned by
\texttt{SYMP-001-COUNTABLE-ASSEMBLY-2026-06-24}.}
\end{theorem}

\begin{proof}
For finite \(\Gamma\), the Pauli label space \(V_\Gamma\) is
finite-dimensional.  Apply \texttt{SYMP-COUNT-FIN-001} to the canonical
symplectic form \(\sigma\) on \(\texttt{PauliLabel}\).
\end{proof}

\begin{theorem}[CountableInfiniteAdaptedSymplecticRecursion, ID: SYMP-INF-REC-001, project proof note 004 Lemma 3.1, PDF pp.~3--4]
Let \(V\) be a countable \(\Ftwo\)-vector space that is not
finite-dimensional.  Let \(B\) be a nondegenerate alternating bilinear form on
\(V\), and let \(W\subseteq V\) be isotropic.  Then there are sequences
\[
  z,x:\mathbb N\longrightarrow V
\]
and a subset \(P\subseteq\mathbb N\) such that
\[
  B(z_i,z_j)=0,\qquad B(x_i,x_j)=0,\qquad
  B(z_i,x_j)=\delta_{ij},
\]
the combined family \(z,x\) is a basis of \(V\), and
\[
  W=\spanF\{z_i:i\in P\}.
\]
\source{Project proof note 004, Lemma 3.1, countable-dimensional recursive
Gram--Schmidt construction, PDF pp.~3--4.}
\depends{\texttt{SYMP-COUNT-BASIS-001},
\texttt{SYMP-FIN-RESTRICT-ISCOMPL-001},
\texttt{SYMP-ORTH-RESTRICT-NONDEG-001}, \texttt{SYMP-LININD-001},
\texttt{SYMP-HYP-EXT-001}.}
\lean{Formalization.Stabilizer.exists_adapted_symplectic_countable_infinite}
\status{blocked false as stated by
\texttt{SYMP-INF-REC-001-2026-06-24}; see
\texttt{SYMP-ADAPTED-CLOSED-001} and \texttt{SYMP-INF-CE-001}.}
\end{theorem}

\begin{proof}
There is no proof of this unrestricted statement.  The Lean-checked theorem
\texttt{SYMP-ADAPTED-CLOSED-001} proves that any adapted symplectic basis with
\(W=\spanF\{z_i:i\in P\}\) forces \(W\) to be \(B\)-closed:
\[
  W^{\perp\perp}=W .
\]
The Lean-checked counterexample \texttt{SYMP-INF-CE-001} gives an isotropic
subspace of \(\texttt{PauliLabel}\,\mathbb N\) that is not \(B\)-closed and
therefore has no adapted symplectic basis.  The residual-step proof sketch
above was missing precisely this closedness hypothesis: without it, the
spanning-side partner need not remain in \(W^\perp\), and the final
\(W=\spanF\{z_i:i\in P\}\) conclusion fails.
\end{proof}

\begin{lemma}[SymplecticBasisOfSpanning, ID: SYMP-SPANNING-BASIS-001]
Let \(B\) be an \(\Ftwo\)-valued bilinear form on \(V\), let \(I\) be an
index type, and let \(z,x:I\to V\) satisfy
\[
  B(z_i,z_j)=0,\qquad B(x_i,x_j)=0,\qquad
  B(z_i,x_j)=\delta_{ij}.
\]
If the combined range of \(z\) and \(x\) spans \(V\), then the combined
family \(I\oplus I\to V\) is a basis of \(V\).
\depends{\texttt{SYMP-LININD-001}.}
\lean{Formalization.Stabilizer.symplectic_basis_of_spanning}
\status{lean_checked support target returned by
\texttt{SYMP-INF-REC-001-2026-06-24}.}
\end{lemma}

\begin{proof}
\texttt{SYMP-LININD-001} gives linear independence from the symplectic
pairing relations.  Combining linear independence with the spanning
hypothesis gives a basis by \(\texttt{Module.Basis.mk}\).
\end{proof}

\begin{theorem}[AdaptedSymplecticBasisImpliesClosed, ID: SYMP-ADAPTED-CLOSED-001]
Let \(B\) be an alternating \(\Ftwo\)-valued bilinear form on \(V\).  Suppose
\(z,x:I\to V\) satisfy the symplectic pairing relations, suppose
\(\{z_i,x_i:i\in I\}\) is a basis of \(V\), and suppose
\[
  W=\spanF\{z_i:i\in P\}
\]
for some \(P\subseteq I\).  Then \(W\) is \(B\)-closed:
\[
  W^{\perp\perp}=W.
\]
\depends{\texttt{SYMP-LININD-001}.}
\lean{Formalization.Stabilizer.adapted_symplectic_closed}
\status{lean_checked obstruction target returned by
\texttt{SYMP-INF-REC-001-2026-06-24}.}
\end{theorem}

\begin{proof}
Expand an arbitrary vector in the adapted symplectic basis.  If
\(u\in W^{\perp\perp}\), then pairing \(u\) against \(x_j\) for
\(j\notin P\) kills the coefficient of \(z_j\), while pairing \(u\) against
each \(z_j\) kills the coefficient of \(x_j\).  Thus \(u\) is a finite
linear combination only of the \(z_i\) with \(i\in P\), hence \(u\in W\).
The reverse inclusion \(W\subseteq W^{\perp\perp}\) follows from the
definition of orthogonal complement.
\end{proof}

\begin{theorem}[InfiniteNonclosedIsotropicCounterexample, ID: SYMP-INF-CE-001]
In the Pauli label space \(\texttt{PauliLabel}\,\mathbb N\), let
\[
  W_{\mathrm{cex}}
  =
  \spanF\{\,Z_i+Z_{i+1}:i\in\mathbb N\,\}.
\]
Then \(W_{\mathrm{cex}}\) is isotropic but not \(B\)-closed.  In particular,
there are no sequences \(z,x:\mathbb N\to\texttt{PauliLabel}\,\mathbb N\)
and no \(P\subseteq\mathbb N\) satisfying the unrestricted
\texttt{SYMP-INF-REC-001} conclusion for \(W_{\mathrm{cex}}\).
\depends{\texttt{PAULI-001}, \texttt{SYMP-ADAPTED-CLOSED-001}.}
\lean{Formalization.Stabilizer.cexW_no_adapted_symplectic_basis}
\status{lean_checked counterexample returned by
\texttt{SYMP-INF-REC-001-2026-06-24}.}
\end{theorem}

\begin{proof}
The generators \(Z_i+Z_{i+1}\) have no \(X\)-coordinate, so their
symplectic pairings with each other vanish and \(W_{\mathrm{cex}}\) is
isotropic.  The vector \(Z_0\) is not in \(W_{\mathrm{cex}}\): the total
\(Z\)-coordinate sum vanishes on each generator but is \(1\) on \(Z_0\).
On the other hand, any vector orthogonal to all generators has equal
successive \(X\)-coordinates; finite support forces all of those
coordinates to vanish, so \(Z_0\in W_{\mathrm{cex}}^{\perp\perp}\).  Thus
\(W_{\mathrm{cex}}^{\perp\perp}\ne W_{\mathrm{cex}}\).  By
\texttt{SYMP-ADAPTED-CLOSED-001}, an adapted symplectic basis for
\(W_{\mathrm{cex}}\) would force closedness, contradiction.
\end{proof}

\begin{lemma}[ClosedOrthogonalPartner, ID: SYMP-CLOSED-PARTNER-001]
Let \(B\) be an alternating \(\Ftwo\)-valued bilinear form on \(V\), let
\(W\subseteq V\) satisfy \(W^{\perp\perp}=W\), and let
\[
  r\in W^\perp\setminus W.
\]
Then there exists \(s\in W^\perp\) such that
\[
  B(r,s)=1.
\]
\depends{\texttt{SYMP-ADAPTED-CLOSED-001}.}
\lean{Formalization.Stabilizer.closed_orthogonal_partner}
\status{proved and Lean-checked from Aristotle task
\texttt{64779911-a022-4fb9-bd7a-42b99cc12d74}.}
\end{lemma}

\begin{proof}
If \(B(r,s)=0\) for every \(s\in W^\perp\), then by definition
\(r\in W^{\perp\perp}\).  The closedness hypothesis gives
\(r\in W\), contradicting \(r\notin W\).  Therefore some \(s\in W^\perp\)
has \(B(r,s)\ne0\); over \(\Ftwo\), this means \(B(r,s)=1\).
\end{proof}

\subsection*{Closed-\(W\) Recursion Support Lemmas}

The following project-local support lemmas are the mathematical spine of the
Lean proof of \texttt{SYMP-INF-REC-CLOSED-001}.  The Lean file also contains
recursion-state records and stability lemmas used to encode the dovetailing;
those are implementation plumbing for this proof, not separate mathematical
targets.

\begin{lemma}[SupHyperbolicRestrictNondegenerate, ID: SYMP-SUP-HYP-NONDEG-001]
Let \(H\subseteq V\) be finite-dimensional and suppose \(B|_H\) is
nondegenerate.  If \(z_0,x_0\in H^\perp\) and \(B(z_0,x_0)=1\), then the
restriction of \(B\) to
\[
  H+\spanF\{z_0,x_0\}
\]
is nondegenerate.
\depends{\texttt{SYMP-FIN-RESTRICT-ISCOMPL-001},
\texttt{SYMP-ORTH-RESTRICT-NONDEG-001}.}
\lean{Formalization.Stabilizer.sup_hyperbolic_restrict_nondegenerate}
\status{proved and Lean-checked from Aristotle task
\texttt{3ae275d6-ba2a-4849-942b-6f2cb283e913}.}
\end{lemma}

\begin{proof}
The summand \(H\) is orthogonal to \(\spanF\{z_0,x_0\}\) because both new
vectors lie in \(H^\perp\).  The restriction on \(H\) is nondegenerate by
hypothesis, and the restriction on \(\spanF\{z_0,x_0\}\) is the two-dimensional
hyperbolic form.  A vector orthogonal to the whole sum has zero \(H\)-component
and zero hyperbolic component, hence is zero.
\end{proof}

\begin{lemma}[OrthogonalComponentMemOrthogonal, ID: SYMP-ORTH-COMP-001]
Assume \(B|_H\) is nondegenerate and
\[
  W=(W\cap H)+(W\cap H^\perp).
\]
If \(t=h+r\in W^\perp\), with \(h\in H\) and \(r\in H^\perp\), then
\[
  h\in W^\perp,\qquad r\in W^\perp .
\]
\depends{\texttt{SYMP-FIN-RESTRICT-ISCOMPL-001}.}
\lean{Formalization.Stabilizer.orthogonal_component_mem_orthogonal}
\status{proved and Lean-checked from Aristotle task
\texttt{3ae275d6-ba2a-4849-942b-6f2cb283e913}.}
\end{lemma}

\begin{proof}
Write every \(w\in W\) as \(w=w_H+w_\perp\) with
\(w_H\in W\cap H\) and \(w_\perp\in W\cap H^\perp\).  Since
\(t\in W^\perp\), \(B(h+r,w)=0\).  The cross pairings
\(B(h,w_\perp)\) and \(B(r,w_H)\) vanish by orthogonality of \(H\) and
\(H^\perp\), so separately \(B(h,w_H)=0\) and \(B(r,w_\perp)=0\).  Varying
\(w\) gives \(h,r\in W^\perp\).
\end{proof}

\begin{lemma}[AdaptedWsplitPreserved, ID: SYMP-WSPLIT-PRESERVE-001]
Assume \(W\) is isotropic and
\[
  W=(W\cap H)+(W\cap H^\perp).
\]
Let \(z_0,x_0\in H^\perp\) satisfy \(B(z_0,x_0)=1\).  If either
\[
  z_0\in W
\]
or
\[
  z_0,x_0\in W^\perp,
\]
then the same split invariant holds after adjoining the pair:
\[
  W=\bigl(W\cap (H+\spanF\{z_0,x_0\})\bigr)
    +\bigl(W\cap (H+\spanF\{z_0,x_0\})^\perp\bigr).
\]
\depends{\texttt{SYMP-SUP-HYP-NONDEG-001},
\texttt{SYMP-ORTH-COMP-001}.}
\lean{Formalization.Stabilizer.adapted_Wsplit_preserved}
\status{proved and Lean-checked from Aristotle task
\texttt{3ae275d6-ba2a-4849-942b-6f2cb283e913}.}
\end{lemma}

\begin{proof}
Use the old split to write \(w\in W\) as \(a+c\).  In the \(P\)-pair case
\(z_0\in W\), decompose the \(H^\perp\)-part into its \(z_0\)-coefficient and
the residual orthogonal to the new hyperbolic plane; isotropy of \(W\) forces
the residual to stay in \(W\).  In the \(O\)-pair case
\(z_0,x_0\in W^\perp\), \(W\) is already orthogonal to the new hyperbolic
plane, so the old \(H^\perp\)-part lies in the new orthogonal complement.
\end{proof}

\begin{lemma}[AdaptedPairFromResidual, ID: SYMP-PAIR-RESIDUAL-001]
Assume \(B\) is nondegenerate and alternating, \(W\) is isotropic and
\(B\)-closed, \(H\) is finite-dimensional with \(B|_H\) nondegenerate, and
\[
  W=(W\cap H)+(W\cap H^\perp).
\]
For every nonzero residual \(r\in H^\perp\), there are
\(z_0,x_0\in H^\perp\) and a tag declaring either a \(P\)-pair or an
\(O\)-pair such that
\[
  B(z_0,x_0)=1,\qquad
  r\in \spanF\{z_0,x_0\},
\]
the new pair preserves nondegeneracy and the \(W\)-split invariant, and either
\(z_0\in W\) or \(z_0,x_0\in W^\perp\) according to the tag.
\depends{\texttt{SYMP-CLOSED-PARTNER-001},
\texttt{SYMP-SUP-HYP-NONDEG-001}, \texttt{SYMP-ORTH-COMP-001},
\texttt{SYMP-WSPLIT-PRESERVE-001}.}
\lean{Formalization.Stabilizer.adapted_pair_from_residual}
\status{proved and Lean-checked from Aristotle task
\texttt{3ae275d6-ba2a-4849-942b-6f2cb283e913}.}
\end{lemma}

\begin{proof}
There are three cases.  If \(r\in W\), take \(z_0=r\) and use
nondegeneracy on \(H^\perp\) to choose \(x_0\in H^\perp\) with
\(B(r,x_0)=1\).  If \(r\in W^\perp\setminus W\), use
\texttt{SYMP-CLOSED-PARTNER-001} to find a partner in \(W^\perp\), then
project that partner to \(H^\perp\) using the finite splitting; the component
lemma keeps it in \(W^\perp\).  If \(r\notin W^\perp\), choose
\(w\in W\) with \(B(r,w)=1\), split \(w\) along \(H\oplus H^\perp\), and use
the \(H^\perp\)-component as the \(W\)-partner.  The previous support lemmas
verify nondegeneracy and preservation of the split invariant in all cases.
\end{proof}

\begin{lemma}[AdaptedSymplecticStepClosed, ID: SYMP-STEP-CLOSED-001]
Under the hypotheses of \texttt{SYMP-PAIR-RESIDUAL-001}, and assuming
\(H^\perp\) contains a nonzero vector, every target \(t\in V\) is covered after
one adapted extension: there are \(z_0,x_0\in H^\perp\) and a \(P/O\)-tag such
that
\[
  t\in H+\spanF\{z_0,x_0\},
\]
the hyperbolic relations hold, and the nondegeneracy and \(W\)-split
invariants are preserved.
\depends{\texttt{SYMP-FIN-RESTRICT-ISCOMPL-001},
\texttt{SYMP-PAIR-RESIDUAL-001}.}
\lean{Formalization.Stabilizer.adapted_symplectic_step}
\status{proved and Lean-checked from Aristotle task
\texttt{3ae275d6-ba2a-4849-942b-6f2cb283e913}.}
\end{lemma}

\begin{proof}
Decompose \(t=h+r\) with \(h\in H\) and \(r\in H^\perp\).  If \(r=0\), use
any nonzero vector in \(H^\perp\) and apply the residual-pair lemma; \(t\) is
already in \(H\).  If \(r\ne0\), apply the residual-pair lemma directly to
\(r\), which puts \(r\), and hence \(t\), in the enlarged span.
\end{proof}

\begin{lemma}[AdaptedWEqSpan, ID: SYMP-W-EQ-SPAN-001]
Let \(z,x:I\to V\) satisfy the symplectic relations and let
\(\operatorname{Sum.elim}(z,x)\) be a basis.  Suppose \(W\) is isotropic,
\(z_i\in W\) for \(i\in P\), and \(z_i,x_i\in W^\perp\) for \(i\notin P\).
Then
\[
  W=\spanF\{z_i:i\in P\}.
\]
\depends{\texttt{SYMP-LININD-001}, \texttt{SYMP-SPANNING-BASIS-001}.}
\lean{Formalization.Stabilizer.adapted_W_eq_span}
\status{proved and Lean-checked from Aristotle task
\texttt{3ae275d6-ba2a-4849-942b-6f2cb283e913}.}
\end{lemma}

\begin{proof}
Expand \(w\in W\) in the symplectic basis.  Pairing \(w\) with each
\(z_i\) kills every \(x_i\)-coefficient, and pairing with \(x_i\) for
\(i\notin P\) kills the \(z_i\)-coefficient outside \(P\).  Thus every
\(w\in W\) lies in the span of \(\{z_i:i\in P\}\).  The reverse inclusion is
the assumption \(z_i\in W\) for \(i\in P\).
\end{proof}

\begin{lemma}[FreshOrthogonalOfInfinite, ID: SYMP-FRESH-PERP-001]
If \(V\) is not finite-dimensional, \(H\subseteq V\) is finite-dimensional,
and \(B|_H\) is nondegenerate, then \(H^\perp\) contains a nonzero vector.
\depends{\texttt{SYMP-FIN-RESTRICT-ISCOMPL-001}.}
\lean{Formalization.Stabilizer.exists_nonzero_mem_orthogonal_of_infinite}
\status{proved and Lean-checked from Aristotle task
\texttt{3ae275d6-ba2a-4849-942b-6f2cb283e913}.}
\end{lemma}

\begin{proof}
If \(H^\perp=\{0\}\), then the complement splitting from
\texttt{SYMP-FIN-RESTRICT-ISCOMPL-001} gives \(V=H\), contradicting that
\(V\) is not finite-dimensional.
\end{proof}

\begin{lemma}[RestrictBotNondegenerate, ID: SYMP-BOT-NONDEG-001]
The restriction of any bilinear form to the zero submodule \(\{0\}\) is
nondegenerate.
\lean{Formalization.Stabilizer.restrict_bot_nondegenerate}
\status{proved and Lean-checked from Aristotle task
\texttt{3ae275d6-ba2a-4849-942b-6f2cb283e913}.}
\end{lemma}

\begin{proof}
The zero submodule has only the zero vector, so both left and right kernels of
the restricted form are trivial.
\end{proof}

\begin{theorem}[ExistsAdaptedSymplecticSequencesClosed, ID: SYMP-SEQUENCES-CLOSED-001]
Under the hypotheses of \texttt{SYMP-INF-REC-CLOSED-001}, there exist
sequences \(z,x:\mathbb N\to V\) and \(P\subseteq\mathbb N\) such that the
symplectic pairing relations hold, \(z_i\in W\) for \(i\in P\),
\(z_i,x_i\in W^\perp\) for \(i\notin P\), and
\[
  \spanF\{z_i,x_i:i\in\mathbb N\}=V.
\]
\depends{\texttt{SYMP-BOT-NONDEG-001}, \texttt{SYMP-FRESH-PERP-001},
\texttt{SYMP-STEP-CLOSED-001}.}
\lean{Formalization.Stabilizer.exists_adapted_symplectic_sequences}
\status{proved and Lean-checked from Aristotle task
\texttt{3ae275d6-ba2a-4849-942b-6f2cb283e913}.}
\end{theorem}

\begin{proof}
Choose a surjective enumeration \(e:\mathbb N\to V\).  Start with
\(H_0=0\), whose restriction is nondegenerate by
\texttt{SYMP-BOT-NONDEG-001}.  At stage \(k\), apply
\texttt{SYMP-STEP-CLOSED-001} to \(e_k\) and the current finite
nondegenerate span \(H_k\).  This produces a new hyperbolic pair and preserves
the symplectic, \(P/O\), nondegeneracy, and \(W\)-split invariants.  Dovetailing
over the enumeration ensures every \(e_k\), hence every vector of \(V\), lies
in the span of the limiting family.
\end{proof}

\begin{theorem}[CountableInfiniteAdaptedSymplecticRecursionClosed, ID: SYMP-INF-REC-CLOSED-001]
Let \(V\) be a countable \(\Ftwo\)-vector space that is not
finite-dimensional.  Let \(B\) be a nondegenerate alternating bilinear form on
\(V\), and let \(W\subseteq V\) be isotropic and \(B\)-closed:
\[
  W^{\perp\perp}=W.
\]
Then there are sequences \(z,x:\mathbb N\to V\) and a subset
\(P\subseteq\mathbb N\) such that the symplectic pairing relations hold, the
combined family \(z,x\) is a basis of \(V\), and
\[
  W=\spanF\{z_i:i\in P\}.
\]
\depends{\texttt{SYMP-COUNT-BASIS-001},
\texttt{SYMP-FIN-RESTRICT-ISCOMPL-001},
\texttt{SYMP-ORTH-RESTRICT-NONDEG-001}, \texttt{SYMP-LININD-001},
\texttt{SYMP-HYP-EXT-001}, \texttt{SYMP-SPANNING-BASIS-001},
\texttt{SYMP-CLOSED-PARTNER-001}, \texttt{SYMP-SUP-HYP-NONDEG-001},
\texttt{SYMP-ORTH-COMP-001}, \texttt{SYMP-WSPLIT-PRESERVE-001},
\texttt{SYMP-PAIR-RESIDUAL-001}, \texttt{SYMP-STEP-CLOSED-001},
\texttt{SYMP-W-EQ-SPAN-001}, \texttt{SYMP-FRESH-PERP-001},
\texttt{SYMP-BOT-NONDEG-001}, \texttt{SYMP-SEQUENCES-CLOSED-001}.}
\lean{Formalization.Stabilizer.exists_adapted_symplectic_countable_infinite_closed}
\status{proved and Lean-checked from Aristotle task
\texttt{3ae275d6-ba2a-4849-942b-6f2cb283e913}.}
\end{theorem}

\begin{proof}
The proof is the dovetailed recursive Gram--Schmidt construction, now with
the missing closedness hypothesis made explicit.  At stages assigned to
basis vectors of \(W\), the residual remains in \(W\) by finite correction
inside the already represented \(W\)-span.  At stages assigned to spanning
directions outside \(W\), the new residual lies in \(W^\perp\setminus W\);
\texttt{SYMP-CLOSED-PARTNER-001} supplies a partner still inside \(W^\perp\).
The finite nondegenerate restriction and orthogonal-complement lemmas preserve
the finite-stage hyperbolic invariant.  The limiting family spans by the
dovetailed enumeration, is linearly independent by \texttt{SYMP-LININD-001},
and becomes a basis by \texttt{SYMP-SPANNING-BASIS-001}.  The \(W\)-stage
invariant gives exactly \(W=\spanF\{z_i:i\in P\}\) by
\texttt{SYMP-W-EQ-SPAN-001}.
\end{proof}

\begin{lemma}[AdaptedSymplecticBasis, ID: SYMP-001, project proof note 004 Lemma 3.1, PDF pp.~3--4]
Let \(V_\Lambda\) be finite or countably dimensional over \(\Ftwo\), with
nondegenerate symplectic form \(\sigma\), and let
\(W\subseteq V_\Lambda\) be isotropic.  There is a symplectic basis
\[
  \{z_i,x_i:i\in I\}
\]
of \(V_\Lambda\), with \(I\) finite or countable, and a subset \(P\subseteq I\)
such that
\[
  W=\spanF\{z_i:i\in P\}.
\]
\source{Project proof note 004, Lemma 3.1, PDF pp.~3--4.}
\depends{\texttt{PAULI-001}, \texttt{SYMP-FIN-001},
\texttt{SYMP-COUNT-BASIS-001},
\texttt{SYMP-FIN-RESTRICT-ISCOMPL-001},
\texttt{SYMP-ORTH-RESTRICT-NONDEG-001}, \texttt{SYMP-LININD-001},
\texttt{SYMP-HYP-EXT-001}, \texttt{SYMP-COUNT-FIN-001},
\texttt{SYMP-PAULI-FIN-COUNT-001}, \texttt{SYMP-INF-REC-CLOSED-001}.}
\lean{Formalization.Stabilizer.adaptedSymplecticBasis}
\status{blueprint_note_complete; Lean checked for the finite subcase and
the listed countable-assembly support lemmas only; the unrestricted
infinite-dimensional recursive assembly is false as stated, and downstream
countable uses must prove \(B\)-closedness or use
\texttt{SYMP-INF-REC-CLOSED-001}.}
\end{lemma}

\begin{proof}
For the finite-dimensional case, use \texttt{SYMP-FIN-001}.  In the countable
case, start with any countable basis \(\{w_i:i\in P\}\) of \(W\).  The
construction is recursive and replaces this basis by another basis
\(\{z_i:i\in P\}\) of the same subspace.  At a finite stage, suppose
hyperbolic pairs have already been chosen with
\(\sigma(z_i,x_j)=\delta_{ij}\) and all \(z\)-\(z\) and \(x\)-\(x\) pairings
zero.  For the next basis vector \(w_k\in W\), first set
\[
  z_k=w_k+\sum_{i<k}\sigma(w_k,x_i)z_i .
\]
This vector still lies in \(W\), is independent modulo the previous
\(z_i\)'s, and satisfies \(\sigma(z_k,x_i)=0\) for every \(i<k\), while
isotropy of \(W\) gives \(\sigma(z_k,z_i)=0\).

Nondegeneracy gives some candidate \(y\) with \(\sigma(z_k,y)=1\).  Replace
\[
  y\quad\text{by}\quad
  y+\sum_{i<k}\sigma(y,x_i)z_i+\sum_{i<k}\sigma(y,z_i)x_i
\]
to make it orthogonal to all previously chosen \(z_i\)'s and \(x_i\)'s.  This
correction preserves \(\sigma(z_k,y)=1\) because \(z_k\) is already orthogonal
to the previous pairs.  Finally replace \(y\) by
\[
  x_k=y+\sigma(y,y)z_k .
\]
Over \(\Ftwo\) the symplectic form is alternating, so this last correction is
usually vacuous, but it records the finite-stage Gram--Schmidt normalization;
the resulting \(x_k\) has the required pairings with all previous pairs and
with \(z_k\).

After all vectors of \(W\) have been paired, the algebraic span \(H\) of the
chosen hyperbolic pairs is nondegenerate: every vector of \(H\) has finite
support in the chosen pairs, and the first nonzero coefficient is detected by
pairing with the corresponding partner.  Hence
\[
  V_\Lambda=H\oplus H^\perp
\]
and \(H^\perp\) is again symplectic.  Choose a finite or countable symplectic
basis of \(H^\perp\) by the same recursive Gram--Schmidt extension procedure
applied to an enumeration of a vector-space basis.  Combining this basis with
the paired basis of \(W\) gives the required adapted symplectic basis.

The Lean formalization currently contains the finite case and the support
targets \texttt{SYMP-COUNT-BASIS-001},
\texttt{SYMP-FIN-RESTRICT-ISCOMPL-001},
\texttt{SYMP-ORTH-RESTRICT-NONDEG-001}, \texttt{SYMP-LININD-001},
\texttt{SYMP-HYP-EXT-001}, \texttt{SYMP-COUNT-FIN-001},
\texttt{SYMP-PAULI-FIN-COUNT-001}, and the corrected closed-\(W\) support
targets through \texttt{SYMP-INF-REC-CLOSED-001}.  The
\texttt{SYMP-INF-REC-001-2026-06-24} response refuted this unrestricted
countable version.  The corrected countable path
\texttt{SYMP-INF-REC-CLOSED-001}, which requires \(W^{\perp\perp}=W\), is
now Lean-checked; downstream countable applications must prove that closedness
hypothesis before using it.
\end{proof}

\subsection*{Adapted Pauli Isomorphism}

\begin{lemma}[AdaptedPauliIso, ID: SYMP-002, project proof note 004 Lemma 3.2, PDF p.~4]
Let \(\{z_i,x_i:i\in I\}\) be a symplectic basis of \(V_\Lambda\).  Let
\(\cA_I\) be the formal qubit quasi-local algebra generated by Pauli pairs
\(\{Z_i,X_i:i\in I\}\).  There is a \(*\)-isomorphism
\[
  \alpha:\cA_I\longrightarrow \cA_\Lambda
\]
which sends the formal Pauli label
\[
  \sum_{i\in F}(a_i z_i+b_i x_i),\qquad F\Subset I,
\]
to the Pauli operator in \(\cA_\Lambda\) with the same binary label, up to the
fixed representative-section phase convention.
\source{Project proof note 004, Lemma 3.2, PDF p.~4;
\path{blueprints/007-pauli-quasilocal-algebra.tex}, targets
\texttt{PAULI-002} and \texttt{QL-001}.}
\depends{\texttt{SYMP-001}, \texttt{PAULI-002}, \texttt{QL-001},
\texttt{QL-002}.}
\lean{Formalization.Stabilizer.adaptedPauliIso}
\status{blueprint_note_complete}
\end{lemma}

\begin{proof}
For a finite subset \(F\subseteq I\), the vectors
\(\{z_i,x_i:i\in F\}\) form a finite symplectic basis of their span.  Choose
the Pauli representatives \(P(z_i)\) and \(P(x_i)\) using the phase convention
fixed in \path{blueprints/007-pauli-quasilocal-algebra.tex}.  The commutation
relations among these representatives are exactly the standard finite-qubit
relations:
\[
  P(z_i)P(x_j)=(-1)^{\delta_{ij}}P(x_j)P(z_i),
\]
with all \(z\)-\(z\) and \(x\)-\(x\) pairs commuting.  Therefore the finite
assignment
\[
  Z_i\mapsto P(z_i),\qquad X_i\mapsto P(x_i),\qquad i\in F,
\]
extends to a \(*\)-isomorphism from
\(\bigotimes_{i\in F}M_2(\mathbb C)\) onto the finite Pauli algebra generated
by those representatives.

If \(F\subseteq F'\), the finite isomorphisms agree with the unital inclusions
because the same representative section is used on the smaller set.  Taking
the directed union and then the quasi-local norm closure gives the
\(*\)-isomorphism \(\alpha:\cA_I\to\cA_\Lambda\).
\end{proof}

\begin{warning}[Algebraic, not geometrically local]
The symplectic basis in \texttt{SYMP-001} is algebraic.  Its vectors may mix
distant physical sites in \(\Lambda\).  Consequently \(\alpha\) is an
algebraic UHF identification and need not be implemented by a finite-depth
Clifford circuit, quasi-local automorphism, gLU, or any geometrically local
normal form.
\source{Project proof note 004, Warning 6.2, PDF p.~7.}
\status{Scope caveat for \texttt{SYMP-001}, \texttt{SYMP-002}, and
\texttt{PROD-001}.}
\end{warning}

\subsection*{Product Form of the Canonical State}

\begin{proposition}[ProductFormCanonicalState, ID: PROD-001, project proof note 004 Proposition 4.1, PDF pp.~4--5]
Let \(S'\) be a consistent stabilizer extension with internal region subspace
\[
  W'_\Lambda=W'\cap V_\Lambda.
\]
Choose an adapted symplectic basis as in \texttt{SYMP-001}, so that
\[
  W'_\Lambda=\spanF\{z_i:i\in P\}.
\]
Transport the restricted canonical state \(\omega_{S'}|_{\cA_\Lambda}\) along
the isomorphism \(\alpha\) of \texttt{SYMP-002}.  If the stabilizer character
has values \(\epsilon_i=\chi(z_i)\in\{\pm1\}\), the transported state on
\(\cA_I\) is
\[
  \bigotimes_{i\in P}\omega_{Z,\epsilon_i}
  \otimes
  \bigotimes_{i\in I\setminus P}\tr_2,
\]
where \(\omega_{Z,\epsilon_i}\) is the one-site \(Z_i=\epsilon_i\) eigenstate
and \(\tr_2\) is the normalized one-site trace.
\source{Project proof note 004, Proposition 4.1, PDF pp.~4--5;
\path{blueprints/002-canonical-state.tex}, target \texttt{STATE-001};
\path{blueprints/008-stabilizer-local-state-foundations.tex}, targets
\texttt{STATE-003} and \texttt{STATE-004}.}
\depends{\texttt{STATE-001}, \texttt{STATE-003}, \texttt{STATE-004},
\texttt{SYMP-001}, \texttt{SYMP-002}.}
\lean{Formalization.Stabilizer.productFormCanonicalState}
\status{blueprint_note_complete}
\end{proposition}

\begin{proof}
Every formal finite Pauli label has the adapted-coordinate form
\[
  v=\sum_{i\in F}(a_i z_i+b_i x_i),
\]
with \(F\Subset I\).  By the region-restriction theorem \texttt{STATE-004},
the restricted canonical state has a nonzero coefficient at \(v\) exactly when
\[
  v\in W'_\Lambda=\spanF\{z_i:i\in P\}.
\]
In adapted coordinates this is equivalent to
\[
  b_i=0\ \text{for all } i,\qquad
  a_i=0\ \text{for all } i\notin P.
\]
When this condition holds, the coefficient is the stabilizer character value
\[
  \chi(v)=\prod_{i\in P}\epsilon_i^{a_i}.
\]

For \(i\in P\), the one-site \(Z_i=\epsilon_i\) eigenstate has Pauli
coefficients \(1\) on the identity, \(\epsilon_i\) on \(Z_i\), and \(0\) on
\(X_i\) and \(X_iZ_i\).  For \(i\notin P\), the normalized trace has
coefficient \(1\) on the identity and \(0\) on the nonidentity Pauli labels.
The tensor product of these one-site states therefore has exactly the same
finite Pauli coefficients as the transported canonical state.

The finite Pauli basis from \texttt{PAULI-002} determines every finite-local
observable by its coefficients.  Hence the two states agree on every finite
local algebra.  By the quasi-local extension uniqueness in \texttt{STATE-003},
they agree on the quasi-local completion.
\end{proof}

\subsection*{Lean Extraction Notes}

\begin{longtable}{@{}p{0.15\linewidth}p{0.34\linewidth}p{0.41\linewidth}@{}}
\toprule
Target & Intended file & Notes \\
\midrule
\endhead
\texttt{SYMP-001} &
\texttt{Formalization/Stabilizer/SymplecticBasis.lean} &
Algebra over \(\texttt{ZMod 2}\); no topology or geometry.  The countable
case can be staged through finite recursive extension lemmas. \\
\texttt{SYMP-FIN-001} &
\texttt{Formalization/Stabilizer/SymplecticBasis.lean} &
Lean-checked finite-dimensional arbitrary-isotropic adapted basis theorem,
used as the finite-stage engine for \texttt{SYMP-001}. \\
\texttt{SYMP-COUNT-BASIS-001} &
\texttt{Formalization/Stabilizer/SymplecticBasisCountable.lean} &
Lean-checked countable basis extraction for subspaces of countable
\(\Ftwo\)-modules. \\
\texttt{SYMP-FIN-RESTRICT-ISCOMPL-001} &
\texttt{Formalization/Stabilizer/SymplecticBasisCountable.lean} &
Lean-checked finite nondegenerate restriction complement splitting, without
assuming the ambient space finite-dimensional. \\
\texttt{SYMP-ORTH-RESTRICT-NONDEG-001} &
\texttt{Formalization/Stabilizer/SymplecticBasisCountable.lean} &
Lean-checked nondegeneracy of the restricted form on the orthogonal
complement. \\
\texttt{SYMP-LININD-001} &
\texttt{Formalization/Stabilizer/SymplecticBasisCountable.lean} &
Lean-checked linear independence of arbitrary-index symplectic families. \\
\texttt{SYMP-HYP-EXT-001} &
\texttt{Formalization/Stabilizer/SymplecticBasisCountable.lean} &
Lean-checked one-step hyperbolic extension over a finite nondegenerate
subspace. \\
\texttt{SYMP-COUNT-FIN-001} &
\texttt{Formalization/Stabilizer/SymplecticBasisCountable.lean} &
Lean-checked countable-index packaging of the finite-dimensional adapted
basis theorem. \\
\texttt{SYMP-PAULI-FIN-COUNT-001} &
\texttt{Formalization/Stabilizer/SymplecticBasisCountable.lean} &
Lean-checked countable-index output shape for finite
\(\texttt{PauliLabel}\) registers. \\
\texttt{SYMP-INF-REC-001} &
\texttt{Formalization/Stabilizer/SymplecticBasisCountable.lean} &
Blocked false as stated: arbitrary isotropic \(W\) need not be
\(B\)-closed. \\
\texttt{SYMP-SPANNING-BASIS-001} &
\texttt{Formalization/Stabilizer/SymplecticBasisCountable.lean} &
Lean-checked basis construction from symplectic relations plus spanning. \\
\texttt{SYMP-ADAPTED-CLOSED-001} &
\texttt{Formalization/Stabilizer/SymplecticBasisCountable.lean} &
Lean-checked necessary closedness obstruction for adapted symplectic bases. \\
\texttt{SYMP-INF-CE-001} &
\texttt{Formalization/Stabilizer/SymplecticBasisCountable.lean} &
Lean-checked nonclosed isotropic counterexample in
\(\texttt{PauliLabel}\,\mathbb N\). \\
\texttt{SYMP-CLOSED-PARTNER-001} &
\texttt{Formalization/Stabilizer/SymplecticBasisCountable.lean} &
Lean-checked dependency: closedness supplies a partner in \(W^\perp\). \\
\texttt{SYMP-SUP-HYP-NONDEG-001},
\texttt{SYMP-ORTH-COMP-001},
\texttt{SYMP-WSPLIT-PRESERVE-001},
\texttt{SYMP-PAIR-RESIDUAL-001},
\texttt{SYMP-STEP-CLOSED-001},
\texttt{SYMP-W-EQ-SPAN-001},
\texttt{SYMP-FRESH-PERP-001},
\texttt{SYMP-BOT-NONDEG-001},
\texttt{SYMP-SEQUENCES-CLOSED-001} &
\texttt{Formalization/Stabilizer/SymplecticBasisCountable.lean} &
Lean-checked mathematical support package returned with
\texttt{SYMP-INF-REC-CLOSED-001}; recursion-state records and stability lemmas
inside the Lean file are implementation plumbing for this package. \\
\texttt{SYMP-INF-REC-CLOSED-001} &
\texttt{Formalization/Stabilizer/SymplecticBasisCountable.lean} &
Corrected infinite-dimensional countable recursive assembly under
\(W^{\perp\perp}=W\); Lean-checked. \\
\texttt{SYMP-002} &
\texttt{Formalization/Stabilizer/AdaptedPauliIso.lean} &
Build from finite Pauli algebra universality and compatible quasi-local
inclusions.  Countable uses now require the closed version of
\texttt{SYMP-001}; keep representative-section phases explicit. \\
\texttt{PROD-001} &
\texttt{Formalization/Stabilizer/ProductForm.lean} &
Prove equality of states by equality of Pauli coefficients on finite local
algebras, then use the \(\texttt{STATE-003}\) extension/uniqueness interface;
do not rely on unrestricted countable \(\texttt{SYMP-001}\). \\
\bottomrule
\end{longtable}

\begin{thebibliography}{9}
\bibitem{ProofNote004}
Project proof note 004,
\emph{Internal Stabilizer Reduction and Projection-State Type Classification}.
Local source:
\path{../runs/2026-06-18-stabilizer-algebra-classification/artifacts/004-internal-stabilizer-reduction-type-classification-R1.pdf}.
\end{thebibliography}

\end{document}
